\documentclass[12pt]{article}
\usepackage[utf8]{inputenc}
\usepackage[T1]{fontenc}
\usepackage{graphicx}
\usepackage{subcaption}
\usepackage{amsmath}
\usepackage{fancyhdr}
\usepackage{geometry}
\geometry{a4paper, margin=1in}
\usepackage{dirtytalk}
\usepackage[english]{babel}
\usepackage{csquotes}
\usepackage{hyperref}
\usepackage{listings}
\usepackage{xcolor}

\definecolor{codebg}{gray}{0.96}
\lstset{
    language=bash,
    basicstyle=\ttfamily\footnotesize,
    backgroundcolor=\color{codebg},
    frame=single, rulecolor=\color{gray!40},
    breaklines=true, columns=fullflexible,
    keepspaces=true, showstringspaces=false, upquote=true
}

\setlength{\headheight}{14.5pt}
\pagestyle{fancy}
\fancyhf{}
\lhead{Operating Systems -- Lab 1}
\rhead{Florea David, FAF-242}
\cfoot{\thepage}

\begin{document}

\begin{titlepage}
\begin{center}
    \IfFileExists{image.png}{\includegraphics[width=0.3\textwidth]{image.png} \\[0.2cm]}{}

    \textbf{MINISTRY OF EDUCATION, CULTURE AND RESEARCH \\ OF THE REPUBLIC OF MOLDOVA} \\[0.3cm]

    \textbf{Technical University of Moldova \\ Faculty of Computers, Informatics and Microelectronics \\ Department of Software and Automation Engineering} \\[2cm]

    \textbf{Florea David, FAF-242}\\[0.5cm]

    \Huge \textbf{Report} \\[0.5cm]

    \large Laboratory work nr. 1 \\[0.5cm]

    \textbf{The OS as a Resource Manager} \\[3cm]

        \begin{flushright}
        \textit{Checked by:} \\
        \textbf{Reițman Patricia}, \textit{University Assistant} \\
        Department of Software and Automation Engineering,
        \\FCIM Faculty, UTM
    \end{flushright}
    \vfill

    Chișinău -- 2026
\end{center}
\end{titlepage}

\newpage

\section*{Setup}
Track A (Linux). Ubuntu 26.04.1 LTS running in VirtualBox, user \texttt{david}.
\begin{lstlisting}[language=bash]
mkdir -p ~/os-lab1 && cd ~/os-lab1
whoami
uname -a
uptime
\end{lstlisting}
\begin{lstlisting}
david
Linux david-VirtualBox 7.0.0-38-generic #38-Ubuntu SMP PREEMPT_DYNAMIC Fri Sep  4 09:10:14 UTC 2026 x86_64 GNU/Linux
 14:13:02 up  1:28,  1 user,  load average: 1.57, 1.67, 1.47
\end{lstlisting}

\section*{Part 1 -- Files and Directories}
\begin{lstlisting}[language=bash]
ls -la /
\end{lstlisting}
\begin{lstlisting}
drwxr-xr-x  20 root root       4096 Oct  5 12:15 .
lrwxrwxrwx   1 root root          7 Apr 20 11:46 bin -> usr/bin
drwxr-xr-x   3 root root       4096 Oct  5 12:25 boot
drwxr-xr-x  19 root root       4160 Oct  5 14:05 dev
drwxr-xr-x 141 root root      12288 Oct  5 12:57 etc
drwxr-xr-x   3 root root       4096 Oct  5 12:18 home
lrwxrwxrwx   1 root root          7 Apr 20 11:46 lib -> usr/lib
dr-xr-xr-x 251 root root          0 Oct  5 12:41 proc
lrwxrwxrwx   1 root root          8 Apr 20 11:46 sbin -> usr/sbin
-rw-------   1 root root 3530555392 Oct  5 12:15 swap.img
dr-xr-xr-x  13 root root          0 Oct  5 12:41 sys
drwxrwxrwt  17 root root        400 Oct  5 14:12 tmp
drwxr-xr-x  13 root root       4096 Oct  5 12:57 usr
drwxr-xr-x  14 root root       4096 Oct  5 12:29 var
...
\end{lstlisting}

\begin{lstlisting}[language=bash]
cd /etc && ls | head -6
\end{lstlisting}
\begin{lstlisting}
ModemManager
NetworkManager
PackageKit
UPower
X11
adduser.conf
\end{lstlisting}

\begin{lstlisting}[language=bash]
cd ~/os-lab1 && mkdir demo && cd demo
echo "hello operating systems" > note.txt
cp note.txt copy.txt
mv copy.txt renamed.txt
ls -l
rm renamed.txt
ls -l
\end{lstlisting}
\begin{lstlisting}
total 8
-rw-rw-r-- 1 david david 24 Oct  5 14:13 note.txt
-rw-rw-r-- 1 david david 24 Oct  5 14:13 renamed.txt
total 4
-rw-rw-r-- 1 david david 24 Oct  5 14:13 note.txt
\end{lstlisting}

\begin{lstlisting}[language=bash]
chmod 600 note.txt && ls -l note.txt
chmod 644 note.txt && ls -l note.txt
id
\end{lstlisting}
\begin{lstlisting}
-rw------- 1 david david 24 Oct  5 14:13 note.txt
-rw-r--r-- 1 david david 24 Oct  5 14:13 note.txt
uid=1000(david) gid=1000(david) groups=1000(david),4(adm),24(cdrom),27(sudo),...
\end{lstlisting}

\subsection*{Observations}
\textbf{(1) Owner and group.} The files I created (\texttt{note.txt}, \texttt{renamed.txt}) are owned by user \textbf{david} and group \textbf{david}. On Ubuntu every user gets a private group with the same name (\texttt{id} shows \texttt{uid=1000(david) gid=1000(david)}).

\textbf{(2) The ten characters after \texttt{chmod 600}.} \texttt{ls -l note.txt} printed \texttt{-rw-------}:
\begin{center}
\begin{tabular}{ll}
\hline
\texttt{-} & file type: regular file (\texttt{d} = directory, \texttt{l} = symlink, \texttt{c}/\texttt{b} = device) \\
\texttt{rw-} & owner (david) may read and write but not execute \\
\texttt{---} & group (david) gets no access \\
\texttt{---} & others get no access \\
\hline
\end{tabular}
\end{center}
So after \texttt{chmod 600} only I can read or modify the file. After \texttt{chmod 644} it became \texttt{-rw-r--r--}: everyone can read it but only I can write it.

\textbf{(3) Two directories under \texttt{/}.}
\begin{itemize}
    \item \texttt{/etc} holds system-wide configuration files (e.g.\ \texttt{adduser.conf}, \texttt{NetworkManager/}, \texttt{X11/}).
    \item \texttt{/home} holds the users' personal home directories (here \texttt{/home/david}).
    \item (Bonus) \texttt{/proc} is not on disk at all. It's a virtual filesystem the kernel generates, with one directory per running process, which is why its size shows as \texttt{0}.
\end{itemize}

\textbf{Interesting output:} on this Ubuntu, \texttt{/bin}, \texttt{/lib} and \texttt{/sbin} are just symlinks into \texttt{/usr} (the \say{merged /usr} layout), and swap is a 3.3\,GB file (\texttt{/swap.img}) rather than a partition.

\section*{Part 2 -- Processes}
\begin{lstlisting}[language=bash]
ps aux | head -4
ps aux | wc -l
top -b -n 1 | head -12
\end{lstlisting}
\begin{lstlisting}
USER         PID %CPU %MEM    VSZ   RSS TTY      STAT START   TIME COMMAND
root           1  0.0  0.3  25920 13600 ?        Ss   12:44   0:03 /usr/lib/systemd/systemd --switched-root --system --deserialize=52 splash
root           2  0.0  0.0      0     0 ?        S    12:44   0:00 [kthreadd]
root           3  0.0  0.0      0     0 ?        S    12:44   0:00 [pool_workqueue_release]
241
top - 14:13:32 up  1:28,  1 user,  load average: 1.30, 1.60, 1.45
Tasks: 240 total,   1 running, 238 sleeping,   0 stopped,   1 zombie
%Cpu(s):  5.9 us,  5.9 sy,  0.0 ni, 82.4 id,  0.0 wa,  0.0 hi,  5.9 si,  0.0 st
MiB Mem :   3398.8 total,    185.6 free,   1872.1 used,   1498.8 buff/cache
MiB Swap:   3367.0 total,   2388.5 free,    978.5 used.   1526.7 avail Mem

    PID USER      PR  NI    VIRT    RES    SHR S  %CPU  %MEM     TIME+ COMMAND
   3602 david     20   0 1865644 101020  72224 S  18.2   2.9   0:28.22 ptyxis
   2594 david     20   0 4053412 247252 103592 S   9.1   7.1   4:02.74 gnome-s+
  11337 david     20   0   13380   5848   3792 R   9.1   0.2   0:00.01 top
      1 root      20   0   25920  13600   9576 S   0.0   0.4   0:03.51 systemd
      2 root      20   0       0      0      0 S   0.0   0.0   0:00.02 kthreadd
\end{lstlisting}

\begin{lstlisting}[language=bash]
sleep 300 &
jobs -l
ps -ef | grep sleep
kill 11339
jobs -l
\end{lstlisting}
\begin{lstlisting}
[1] 11339
[1]+ 11339 Running                    sleep 300 &
david      11339   11331  0 14:13 ?        00:00:00 sleep 300
david      11341   11331  0 14:13 ?        00:00:00 grep sleep
[1]+ 11339 Terminated                 sleep 300
\end{lstlisting}

\begin{lstlisting}[language=bash]
sleep 300 & PID=$!
ls /proc/$PID/ | head -5
cat /proc/$PID/status | head -6
kill $PID
ps -p 1 -o pid,user,comm,args
ls -l /sbin/init
\end{lstlisting}
\begin{lstlisting}
arch_status
attr
autogroup
auxv
cgroup
Name:	sleep
Umask:	0002
State:	S (sleeping)
Tgid:	11344
Ngid:	0
Pid:	11344
    PID USER     COMMAND         COMMAND
      1 root     systemd         /usr/lib/systemd/systemd --switched-root --system --deserialize=52 splash
lrwxrwxrwx 1 root root 22 Jul 27 23:44 /sbin/init -> ../lib/systemd/systemd
\end{lstlisting}

\subsection*{Observations}
\textbf{(1) PID 1.} PID 1 is \textbf{systemd} (run by root; \texttt{/sbin/init} is a symlink to it). It's the first user-space process the kernel starts at boot (the \emph{init} process). It starts and supervises all the system services, and it adopts orphaned processes. PID 2 is \texttt{kthreadd}, the parent of all kernel threads.

\textbf{(2) Number of processes.} \texttt{ps aux | wc -l} returned \textbf{241}. Subtracting the header line leaves \textbf{240 processes} on an \say{idle} VM, which matches \texttt{top}: \texttt{Tasks: 240 total, 1 running, 238 sleeping, 1 zombie}. Most of them are kernel threads (\texttt{[kworker/...]}, \texttt{[rcu\_...]}) and GNOME desktop processes.

\textbf{(3) State of the sleeping process.} \texttt{/proc/<pid>/status} showed \texttt{State: S (sleeping)}. \texttt{sleep 300} isn't using the CPU. It's blocked in the kernel waiting for a timer to expire, so the scheduler doesn't run it until then.

\textbf{Interesting output:} \texttt{kill} doesn't kill anything by itself. It asks the kernel to deliver signal 15 (SIGTERM) to the process, and the kernel then terminates it, which \texttt{jobs -l} reports as \texttt{Terminated}.

\section*{Part 3 -- Memory}
\begin{lstlisting}[language=bash]
free -h
cat /proc/meminfo | head -6
swapon --show
\end{lstlisting}
\begin{lstlisting}
               total        used        free      shared  buff/cache   available
Mem:           3.3Gi       1.8Gi       192Mi        39Mi       1.5Gi       1.5Gi
Swap:          3.3Gi       978Mi       2.3Gi
MemTotal:        3480328 kB
MemFree:          197556 kB
MemAvailable:    1576952 kB
Buffers:            9060 kB
Cached:          1457600 kB
SwapCached:       104292 kB
NAME      TYPE SIZE   USED PRIO
/swap.img file 3.3G 978.4M   -1
\end{lstlisting}

\begin{lstlisting}[language=bash]
sleep 300 & PID=$!
grep VmRSS /proc/$PID/status
grep -E "VmSize|RssAnon|RssFile|RssShmem" /proc/$PID/status
kill $PID
\end{lstlisting}
\begin{lstlisting}
VmRSS:	    7740 kB
VmSize:	   16116 kB
RssAnon:	    1224 kB
RssFile:	    6516 kB
RssShmem:	       0 kB
\end{lstlisting}

\subsection*{Observations}
\textbf{(1) Total and free RAM.} \texttt{free -h}: \textbf{3.3\,GiB total} (MemTotal 3\,480\,328\,kB), 1.8\,GiB used, \textbf{192\,MiB \say{free}} but \textbf{1.5\,GiB \say{available}}. \say{Free} is low because Linux uses otherwise idle RAM as page cache (\texttt{buff/cache} 1.5\,GiB). That cache is released when programs need memory, so \say{available} is the real headroom.

\textbf{(2) Swap.} Swap is disk space the kernel uses as an overflow for RAM: rarely used memory pages are moved out to disk to make room in physical memory. This VM has a \textbf{3.3\,GiB swap file} (\texttt{/swap.img}, not a partition), of which about 978\,MiB was in use. The VM is fairly memory-constrained with a desktop running.

\textbf{(3) VmRSS of \texttt{sleep}.} \texttt{VmRSS: 7740 kB} (about 7.6\,MB), while the virtual size was \texttt{VmSize: 16116 kB}. That surprised me for a program that does nothing. The breakdown explains it: \texttt{RssAnon: 1224 kB} is its own private memory and \texttt{RssFile: 6516 kB} is mapped files, mostly the C library and the \texttt{sleep} binary. Those file pages are \textbf{shared} with every other process using libc, so \texttt{sleep} only costs about 1.2\,MB of extra RAM.

\section*{Part 4 -- Devices and Storage}
\begin{lstlisting}[language=bash]
df -h
lsblk
findmnt /
\end{lstlisting}
\begin{lstlisting}
Filesystem      Size  Used Avail Use% Mounted on
tmpfs           680M  2.4M  678M   1% /run
/dev/sda2        20G   12G  7.5G  60% /
tmpfs           1.7G  4.0K  1.7G   1% /dev/shm
tmpfs           1.7G  5.0M  1.7G   1% /tmp
tmpfs           340M  5.4M  335M   2% /run/user/1000
/dev/sr0         51M   51M     0 100% /run/media/david/VBox_GAs_7.2.20
NAME   MAJ:MIN RM   SIZE RO TYPE MOUNTPOINTS
sda      8:0    0    20G  0 disk
|-sda1   8:1    0     1M  0 part
`-sda2   8:2    0    20G  0 part /
sr0     11:0    1  50.7M  1 rom  /run/media/david/VBox_GAs_7.2.20
TARGET SOURCE    FSTYPE OPTIONS
/      /dev/sda2 ext4   rw,relatime
\end{lstlisting}
(Output shortened: the two \texttt{none} credential mounts from \texttt{df} and the 13 \texttt{loop} devices from \texttt{lsblk}, which belong to snap packages, are omitted.)

\begin{lstlisting}[language=bash]
ls -l /dev | head -6
ls -l /dev/null /dev/sda /dev/tty /dev/urandom
\end{lstlisting}
\begin{lstlisting}
total 0
crw-r--r--  1 root  root     10, 235 Oct  5 12:41 autofs
drwxr-xr-x  2 root  root        400 Oct  5 13:49 block
drwxr-xr-x  2 root  root         80 Oct  5 12:41 bsg
crw-------  1 root  root     10, 234 Oct  5 12:41 btrfs-control
drwxr-xr-x  3 root  root         60 Oct  5 13:40 bus
crw-rw-rw- 1 root root 1, 3 Oct  5 12:41 /dev/null
brw-rw---- 1 root disk 8, 0 Oct  5 12:41 /dev/sda
crw-rw-rw- 1 root tty  5, 0 Oct  5 12:57 /dev/tty
crw-rw-rw- 1 root root 1, 9 Oct  5 12:41 /dev/urandom
\end{lstlisting}

\subsection*{Observations}
\textbf{(1) Root filesystem device.} \texttt{/} is mounted from \textbf{\texttt{/dev/sda2}}, an \textbf{ext4} partition of 20\,GB (12\,GB used, 60\%). \texttt{lsblk} shows that \texttt{sda} is the VM's single 20\,GB virtual disk, with \texttt{sda1} a tiny 1\,MB BIOS-boot partition and \texttt{sda2} the root.

\textbf{(2) A \texttt{/dev} entry.}
\begin{itemize}
    \item \texttt{/dev/sda} (\texttt{brw-rw---- root disk 8, 0}) is the whole virtual hard disk: \texttt{b} = block device, major number 8 = SCSI/SATA disk driver.
    \item \texttt{/dev/null} (\texttt{crw-rw-rw- 1, 3}) is a character device that discards anything written to it.
    \item \texttt{/dev/sr0} is the virtual CD drive, currently holding the VirtualBox Guest Additions ISO.
\end{itemize}

\textbf{(3) \say{Everything is a file}.} Disks, terminals, the CD-ROM, even random numbers (\texttt{/dev/urandom}) and process information (\texttt{/proc}) are all exposed as entries in the same directory tree. Programs use them through the same \texttt{open/read/write/close} calls they use for ordinary files, and the kernel's drivers do the device-specific work behind that interface.

\textbf{Interesting output:} \texttt{df -h} shows that several \say{filesystems} don't live on any disk at all. They are RAM-backed \texttt{tmpfs} mounts (\texttt{/run}, \texttt{/tmp}, \texttt{/dev/shm}).

\section*{Optional -- htop and strace}
\begin{lstlisting}[language=bash]
htop        # captured via a detached GNU screen session
\end{lstlisting}
\begin{lstlisting}
    0[||||                       10.4%] Tasks: 148, 642 thr, 89 kthr; 1 runnin
    1[|||                         6.8%] Load average: 0.89 1.44 1.41
  Mem[|||||||||||||||||||||1.69G/3.32G] Uptime: 01:29:53
  Swp[||||||||||            877M/3.29G]
\end{lstlisting}

\begin{lstlisting}[language=bash]
strace -c ls
\end{lstlisting}
\begin{lstlisting}
demo
% time     seconds  usecs/call     calls    errors syscall
------ ----------- ----------- --------- --------- ------------------
 30.18    0.000497         497         1           execve
 26.29    0.000433          16        27           mmap
 18.52    0.000305         305         1           write
  5.28    0.000087           7        11         1 openat
  4.19    0.000069           7         9           fstat
  3.95    0.000065           6        10           close
  3.76    0.000062           4        14           read
  2.73    0.000045           9         5           brk
  2.00    0.000033           4         8           mprotect
  ...
  0.00    0.000000           0         7         7 ioctl
  0.00    0.000000           0         2           getdents64
------ ----------- ----------- --------- --------- ------------------
100.00    0.001647          11       138        12 total
\end{lstlisting}

\texttt{htop 3.4.1} shows the same data as \texttt{top}, but with per-CPU bars (2 vCPUs here), memory/swap bars, and a scrollable process list sorted by CPU.

\texttt{strace -c ls} made \textbf{138 system calls (12 of them failed)} just to print one word:
\begin{itemize}
    \item \texttt{mmap} $\times$27, \texttt{openat} $\times$11, \texttt{read} $\times$14 and \texttt{mprotect} $\times$8 load the program and its shared libraries into memory;
    \item \texttt{getdents64} $\times$2 is the call that actually reads the directory entries;
    \item \texttt{write} $\times$1 prints \texttt{demo} to stdout;
    \item the 7 failed \texttt{ioctl} calls are \texttt{ls} checking whether stdout is a terminal (it wasn't: \texttt{ENOTTY}).
\end{itemize}
So almost all of the work is setup. The \say{real} job of \texttt{ls} comes down to \texttt{openat} + \texttt{getdents64} + \texttt{write}, and every one of those is a request to the kernel.

\break

\section*{Conclusion}
In this laboratory work I looked at Linux from the outside, as a manager of four kinds of resources. It manages \textbf{files}: \texttt{ls -l} and \texttt{chmod} showed it naming data in one directory tree and enforcing who may read or write each file. It manages \textbf{processes}: \texttt{ps aux} and \texttt{top} showed about 240 processes that it creates, schedules and signals (\texttt{kill}), all visible under \texttt{/proc}. It manages \textbf{memory}: \texttt{free -h} and \texttt{VmRSS} showed it splitting 3.3\,GiB of RAM plus swap between processes and sharing library pages between them. It manages \textbf{devices}: \texttt{df -h}, \texttt{lsblk} and \texttt{ls -l /dev} showed it hiding the disk and other hardware behind file-like interfaces. Finally, \texttt{strace} showed that even a trivial command like \texttt{ls} gets everything it needs from the kernel through system calls.

\end{document}
